\chapter{Presupuesto y viabilidad económica}

\section{Costes}

Para una estimación simple y coherente, se ha considerado un periodo de trabajo de febrero a junio (aprox. 21 semanas) con una dedicación media de 20 horas semanales del estudiante.

\textbf{Hipótesis utilizadas:}
\begin{itemize}
	\item Horas de estudiante: $21 \times 20 = 420$ h.
	\item Coste hora estudiante (ingeniero pre-júnior): 18 EUR/h.
	\item Dedicación de dirección académica (profesor): 40 h totales.
	\item Coste hora profesor (UPC Terrassa): 70 EUR/h.
	\item Google Colab Pro: 100 unidades informáticas con coste total de 10 EUR.
\end{itemize}

\begin{table}[htbp]
\centering
\caption{Resumen de costes del TFG (estimación simple)}
\begingroup
\small
\setlength{\tabcolsep}{4pt}
\begin{tabular}{p{5.2cm}rrr}
\toprule
\textbf{Concepto} & \textbf{Horas / unid.} & \textbf{Coste unitario} & \textbf{Coste total} \\
\midrule
Estudiante (ingeniero pre-júnior) & 420 h & 18 EUR/h & 7.560 EUR \\
Profesor (dirección y seguimiento) & 40 h & 70 EUR/h & 2.800 EUR \\
Google Colab Pro & 100 unidades informáticas & 0,10 EUR/unidad & 10 EUR \\
\midrule
\textbf{Total estimado} &  &  & \textbf{10.370 EUR} \\
\bottomrule
\end{tabular}
\endgroup
\end{table}

Esta estimación debe considerarse orientativa y sin impuestos indirectos, y resulta útil para valorar la viabilidad económica del desarrollo en un contexto académico.


